% The board layer as a contract, and the generator that keeps it honest.
\begin{tikzpicture}[font=\sffamily\scriptsize, x=1cm, y=1cm]

\newcommand{\umlbox}[4]{%
  \node[umlclass, text width=36mm] (#1) at (#2) {%
    \begin{minipage}{36mm}
      \centering\bfseries #3\par\vspace{1pt}
      \hrule\vspace{2pt}
      \raggedright\normalfont #4\vspace{1pt}
    \end{minipage}};}

% ---------------- the description and its generator
\umlbox{desc}{-5.0,3.4}{nucleo\_h7a3zi\_q.yaml}{%
  pins, clocks, absences\\ one evidence field per row\\[1pt]
  {\itshape the source of truth}}
\umlbox{gen}{-5.0,0.9}{gen\_board.py}{%
  emits three headers\\ refuses any row whose\\ evidence is unconfirmed}

% ---------------- the generated headers
\umlbox{pins}{0.0,3.4}{board\_pins.h}{generated, never edited}
\umlbox{clocks}{0.0,1.7}{board\_clocks.h}{generated, never edited}
\umlbox{caps}{0.0,0.0}{board\_caps.h}{%
  {\tiny BOARD\_HAS\_ETHERNET 0}\\ {\tiny BOARD\_HAS\_CAN\_TRANSCEIVER 0}\\[1pt]
  {\itshape absence, with proof}}

% ---------------- the contract
\umlbox{board}{5.2,1.7}{board.h}{%
  board\_init()\\ board\_led(), board\_button()\\ board\_console\_putc()\\
  board\_clock\_hz()\\[1pt] {\itshape no vendor type in any signature}}

% ---------------- what implements it, and what uses it
\umlbox{impl}{5.2,-1.6}{src/bsp/*.c}{%
  board.c clock.c gpio.c uart.c\\ startup.c system.c\\[1pt]
  {\itshape the only files in the tree that include a vendor header}}
\umlbox{above}{5.2,5.3}{everything above}{%
  control, time, sense, act,\\ bus, mw, safety, update}

\draw[flow] (desc) -- (gen);
\draw[flow] (gen) -- (pins);
\draw[flow] (gen) -- (clocks);
\draw[flow] (gen) -- (caps);
\draw[flow] (pins) -- (board);
\draw[flow] (clocks) -- (board);
\draw[flow] (caps) -- (board);
\draw[flow] (board) -- (above);
\draw[flow] (impl) -- node[lbl, right]{implements} (board);

% the rule, checked rather than trusted
\draw[line width=1.1pt, draw=red!60!black] (2.9,3.7) -- (7.6,3.7);
\node[font=\sffamily\tiny, text=red!60!black, fill=white, inner sep=1.5pt]
  at (5.2,3.7) {check\_layering.py: zero vendor includes above this line};

\node[umlnote, text width=46mm, anchor=north west] at (-7.2,-1.0)
  {A pin number that appears in two places will eventually appear in two
   different forms. Writing the board as data and generating the header removes
   that possibility, and it makes the documentation and the code the same
   artefact rather than two that agree for a while.};
\node[umlnote, text width=46mm, anchor=north west] at (-7.2,-4.0)
  {The layering rule is the reason chapter 20 can build most of this node on a
   host machine with no hardware at all. It is trivial to keep from chapter 4
   onward and expensive to retrofit at chapter 19.};
\end{tikzpicture}
